% Options for packages loaded elsewhere
\PassOptionsToPackage{unicode}{hyperref}
\PassOptionsToPackage{hyphens}{url}
\documentclass[
  indonesian,
  11pt,
]{article}
\usepackage{xcolor}
\usepackage[a4paper,top=22mm,bottom=22mm,left=22mm,right=18mm]{geometry}
\usepackage{amsmath,amssymb}
\setcounter{secnumdepth}{-\maxdimen} % remove section numbering
\usepackage{iftex}
\ifPDFTeX
  \usepackage[T1]{fontenc}
  \usepackage[utf8]{inputenc}
  \usepackage{textcomp} % provide euro and other symbols
\else % if luatex or xetex
  \usepackage{unicode-math} % this also loads fontspec
  \defaultfontfeatures{Scale=MatchLowercase}
  \defaultfontfeatures[\rmfamily]{Ligatures=TeX,Scale=1}
\fi
\usepackage{lmodern}
\ifPDFTeX\else
  % xetex/luatex font selection
\fi
% Use upquote if available, for straight quotes in verbatim environments
\IfFileExists{upquote.sty}{\usepackage{upquote}}{}
\IfFileExists{microtype.sty}{% use microtype if available
  \usepackage[]{microtype}
  \UseMicrotypeSet[protrusion]{basicmath} % disable protrusion for tt fonts
}{}
\makeatletter
\@ifundefined{KOMAClassName}{% if non-KOMA class
  \IfFileExists{parskip.sty}{%
    \usepackage{parskip}
  }{% else
    \setlength{\parindent}{0pt}
    \setlength{\parskip}{6pt plus 2pt minus 1pt}}
}{% if KOMA class
  \KOMAoptions{parskip=half}}
\makeatother
\usepackage{longtable,booktabs,array}
\newcounter{none} % for unnumbered tables
\usepackage{calc} % for calculating minipage widths
% Correct order of tables after \paragraph or \subparagraph
\usepackage{etoolbox}
\makeatletter
\patchcmd\longtable{\par}{\if@noskipsec\mbox{}\fi\par}{}{}
\makeatother
% Allow footnotes in longtable head/foot
\IfFileExists{footnotehyper.sty}{\usepackage{footnotehyper}}{\usepackage{footnote}}
\makesavenoteenv{longtable}
\ifLuaTeX
\usepackage[bidi=basic,shorthands=off,provide=*]{babel}
\else
\usepackage[bidi=default,shorthands=off,provide=*]{babel}
\fi
\ifLuaTeX
  \usepackage{selnolig} % disable illegal ligatures
\fi
\setlength{\emergencystretch}{3em} % prevent overfull lines
\providecommand{\tightlist}{%
  \setlength{\itemsep}{0pt}\setlength{\parskip}{0pt}}
\usepackage{bookmark}
\IfFileExists{xurl.sty}{\usepackage{xurl}}{} % add URL line breaks if available
\urlstyle{same}
\usepackage{fancyhdr}
\usepackage{titlesec}
\usepackage{enumitem}
\usepackage{needspace}
\usepackage{ragged2e}
\usepackage{setspace}
\usepackage{newunicodechar}
\newunicodechar{∅}{\ensuremath{\varnothing}}

\definecolor{Ink}{HTML}{222222}
\definecolor{RuleGray}{HTML}{777777}

\setcounter{secnumdepth}{-1}
\setstretch{1.08}
\setlength{\parindent}{0pt}
\setlength{\parskip}{0.35em}
\setlength{\emergencystretch}{3em}
\raggedbottom
\clubpenalty=10000
\widowpenalty=10000
\displaywidowpenalty=10000

\setlist[itemize]{leftmargin=1.75em,itemsep=0.22em,topsep=0.35em,parsep=0pt}
\setlist[enumerate]{leftmargin=1.9em,itemsep=0.22em,topsep=0.35em,parsep=0pt}
\renewcommand{\tightlist}{%
  \setlength{\itemsep}{0.22em}%
  \setlength{\parskip}{0pt}%
}

\titleformat{\section}
  {\normalfont\Large\bfseries\color{Ink}}
  {}{0pt}{}
\titlespacing*{\section}{0pt}{1.4em}{0.55em}
\titleformat{\subsection}
  {\normalfont\large\bfseries\color{Ink}}
  {}{0pt}{}
\titlespacing*{\subsection}{0pt}{1.1em}{0.4em}
\titleformat{\subsubsection}
  {\normalfont\normalsize\bfseries\color{Ink}}
  {}{0pt}{}
\titlespacing*{\subsubsection}{0pt}{0.9em}{0.3em}

\pagestyle{fancy}
\fancyhf{}
\fancyhead[L]{\footnotesize\textit{Roadmap Skripsi TGA}}
\fancyhead[R]{\footnotesize\textit{Checklist kerja}}
\fancyfoot[C]{\thepage}
\renewcommand{\headrulewidth}{0.3pt}
\renewcommand{\footrulewidth}{0pt}
\fancypagestyle{plain}{%
  \fancyhf{}
  \fancyfoot[C]{\thepage}
  \renewcommand{\headrulewidth}{0pt}
}

\AtBeginEnvironment{quote}{%
  \small
  \setlength{\leftskip}{1.2em}
  \setlength{\rightskip}{1.2em}
}
\AtBeginEnvironment{longtable}{%
  \footnotesize
  \setlength{\tabcolsep}{3pt}
  \renewcommand{\arraystretch}{1.15}
  \sloppy
}
\setlength{\LTpre}{0.5em}
\setlength{\LTpost}{0.8em}
\renewcommand{\arraystretch}{1.18}

% fallback for those not using the hyperref driver hyperxmp:
\makeatletter
\@ifundefined{xmpquote}{\newcommand{\xmpquote}[1]{#1}}{}
\makeatother
\hypersetup{
  pdftitle={Roadmap Skripsi TGA - Borrowed Size dan Agglomeration Shadow Jakarta},
  pdflang={id-ID},
  hidelinks,
  pdfcreator={LaTeX via pandoc}}

\title{Roadmap Skripsi TGA\\[0.35em]
  \large Borrowed Size dan Agglomeration Shadow Jakarta}
\author{}
\date{22 September 2026}


\providecommand{\passthrough}[1]{#1}

\providecommand{\passthrough}[1]{#1}

\providecommand{\passthrough}[1]{#1}

\providecommand{\passthrough}[1]{#1}
\begin{document}
\maketitle
\renewcommand*\contentsname{Daftar Isi}
\setcounter{tocdepth}{2}
\tableofcontents\clearpage
\begin{quote}
\textbf{Status:} checklist kerja untuk mengendalikan penelitian,
penulisan, dan persiapan TGA. Berkas ini bukan naskah skripsi final.

\textbf{Diperbarui:} 22 September 2026

\textbf{Aturan utama:} kotak hanya dicentang setelah pekerjaan
menghasilkan artefak atau keputusan yang dapat diperiksa. Data yang
direncanakan tidak boleh ditulis sebagai data yang sudah diterima atau
valid.
\end{quote}

\section{1. Sumber dan status awal}\label{sumber-dan-status-awal}

Gunakan sumber berikut sebagai dasar kerja, dengan tetap mengikuti
arahan terbaru dosen pembimbing dan kanal resmi TGA:

\begin{itemize}
\tightlist
\item
  Bab 1 Pendahuluan final
\item
  Bab 2 Kajian Pustaka final
\item
  Bab 3 Metode Penelitian final
\item
  Bab 4 Gambaran Umum Lokasi Penelitian final
\item
  {[}{[}source/official-document/Panduan\_TGA{[}1{]}.pdf\textbar Panduan
  TGA{]}{]}
\item
  Template TGA Penelitian 2026
\item
  Skripsi S1 - Borrowed Size dan Agglomeration Shadow Jakarta
\item
  kajian penelitian terdahulu borrowed size dan agglomeration shadow
  jakarta bahan bab 2 2
\end{itemize}

Panduan yang tersedia memuat rujukan tahun 2023, sedangkan template
penelitian bertahun 2026. Persyaratan administratif, formulir, jadwal,
dan mekanisme ujian harus diperiksa ulang pada laman TGA sebelum
pendaftaran.

\subsection{1.1 Identitas kerja}\label{identitas-kerja}

{\def\LTcaptype{none} % do not increment counter
\begin{longtable}[]{@{}
  >{\raggedright\arraybackslash}p{(\linewidth - 2\tabcolsep) * \real{0.5000}}
  >{\raggedright\arraybackslash}p{(\linewidth - 2\tabcolsep) * \real{0.5000}}@{}}
\toprule\noalign{}
\begin{minipage}[b]{\linewidth}\raggedright
Item
\end{minipage} & \begin{minipage}[b]{\linewidth}\raggedright
Isi
\end{minipage} \\
\midrule\noalign{}
\endhead
\bottomrule\noalign{}
\endlastfoot
Nama mahasiswa & {[}isi{]} \\
NIM & {[}isi{]} \\
Dosen pembimbing TGA & {[}isi{]} \\
Judul kerja & Borrowed Size atau Agglomeration Shadow? Analisis Hubungan
Integrasi Metropolitan dengan Kematangan Pusat-Pusat Sekunder di Kawasan
Aglomerasi Jakarta \\
Target periode TGA & {[}isi setelah disepakati dengan dosen
pembimbing{]} \\
Target pra-sidang & {[}isi{]} \\
Target sidang & {[}isi{]} \\
Target yudisium & {[}isi{]} \\
\end{longtable}
}

\subsection{1.2 Posisi penelitian yang sedang
dipertahankan}\label{posisi-penelitian-yang-sedang-dipertahankan}

Penelitian ini sepenuhnya bertumpu pada proksi dan sumber terbuka, tanpa
ketergantungan pada permohonan data tertutup. Data pemerintah yang
terbuka tetap dapat digunakan. Pemilihan sumber mengikuti kesesuaian
konstruk, cakupan spasial dan temporal, keterulangan pengolahan, serta
biaya pemerolehan; status resmi tidak menjadi ukuran mutu dengan
sendirinya.

Desain penelitian bersifat multitemporal sesuai bukti: perubahan
dianalisis pada indikator dengan seri yang sebanding, sedangkan
indikator lain dianalisis pada periode yang tersedia. Tahun dasar dan
rentang analisis ditetapkan per modul setelah audit; panel lengkap dan
satu tahun jangkar untuk semua sumber tidak diasumsikan.

\begin{itemize}
\tightlist
\item
  Kasus utama adalah Jabodetabek/Kawasan Aglomerasi Jakarta.
\item
  DKI Jakarta diperlakukan sebagai satu inti untuk hubungan eksternal,
  dengan batas administratif tetap.
\item
  Tolok ukur utama menggunakan pusat sekunder Jabodetabek yang
  sebanding; DKI tidak membentuk garis ekspektasi. WSM di luar cakupan
  hanya menjadi sensitivitas jika kompatibel.
\item
  Tahun dan rentang analisis ditetapkan per modul setelah audit
  kesebandingan; setiap data tetap memuat periode observasinya.
\item
  Unit utama ditargetkan pada kecamatan hanya jika bukti kritis
  kompatibel. Kabupaten/kota tetap menjadi konteks kebijakan dan
  tanggung jawab pemerintahan.
\item
  Titik dan raster dipakai pada resolusi asalnya sebagai lapisan
  pendukung. Nilai agregat tidak diturunkan secara semu menjadi variasi
  kecamatan, grid, atau titik.
\item
  Proksi, estimasi, dan data pemerintah terbuka dinilai melalui
  kesesuaian konstruk dan validasi, bukan hierarki status lembaga.
\item
  Tipologi deskriptif boleh dibangun sesuai data. Label ketat
  \emph{borrowed size} atau \emph{agglomeration shadow} hanya digunakan
  setelah empat gerbang bukti terpenuhi.
\item
  Web GIS interaktif adalah keluaran tambahan setelah analisis minimum
  selesai. Keluaran ini dapat ditiadakan tanpa mengubah jawaban
  penelitian.
\end{itemize}

\subsection{1.3 Rancangan minimum yang
dikunci}\label{rancangan-minimum-yang-dikunci}

Rantai utama penelitian dibatasi pada:

\begin{enumerate}
\def\labelenumi{\arabic{enumi}.}
\tightlist
\item
  Satu periode yang kompatibel;
\item
  Satu bukti relasi komuter terarah;
\item
  Satu domain fungsi yang dibentuk secara independen, dengan pekerjaan
  sebagai pilihan utama;
\item
  Tolok ukur internal Jabodetabek;
\item
  Residu fungsi relatif dan hubungannya dengan profil relasi; serta
\item
  Diagnosis hanya jika empat gerbang bukti terpenuhi.
\end{enumerate}

Domain kedua, analisis multitemporal, CCTV, SUMO, pembanding eksternal,
dan Web GIS merupakan penguatan opsional. Jika jumlah pusat efektif
tidak mendukung regresi, gunakan profil dan perbandingan kelompok.

\subsection{1.4 Tesis kerja}\label{tesis-kerja}

\begin{quote}
Penelitian ini mengukur apakah integrasi dengan Jakarta memungkinkan
pusat sekunder meminjam manfaat skala metropolitan sambil mematangkan
fungsi lokalnya, atau justru menciptakan ketergantungan asimetris;
kemudian menguji bagaimana jejak kondisi tersebut berubah di bawah
berbagai asumsi indikator, tanpa mengubah batas dan tanggung jawab
administratif.
\end{quote}

Tesis ini tetap merupakan arah kerja sampai dikonfirmasi dalam
pembimbingan. Ia tidak boleh diubah menjadi klaim bahwa Jakarta secara
kausal menyebabkan \emph{upgrading}, ketergantungan, atau \emph{shadow}
tanpa desain identifikasi tambahan.

\section{2. Urutan kendali penelitian}\label{urutan-kendali-penelitian}

Urutan ini menjadi jalur utama. Jika satu gerbang belum valid, kembali
ke tahap sebelumnya; jangan melompati gerbang hanya untuk mengejar
penulisan bab.

\begin{enumerate}
\def\labelenumi{\arabic{enumi}.}
\tightlist
\item
  Kesiapan administratif dan pembimbingan.
\item
  Penguncian judul, pertanyaan, ruang lingkup, dan batas klaim.
\item
  Akuisisi sumber terbuka, konstruksi proksi, dan perancangan
  pemeriksaan independen.
\item
  Audit mutu bukti, harmonisasi, dan pembekuan baseline.
\item
  Validitas data dan Progres 1.
\item
  Delineasi pusat sekunder, wilayah tangkapan, dan unit perhitungan.
\item
  Pengukuran relasi komuter, fungsi relatif pada maksimal dua domain,
  manfaat, beban, dan diagnosis bersyarat.
\item
  Sensitivitas, validitas hasil, dan Progres 2.
\item
  Penulisan lengkap Bab 1-6 dan lampiran.
\item
  Pra-sidang, revisi, sidang, yudisium, dan wisuda.
\end{enumerate}

\section{3. Arahan resmi TGA dan gerbang dosen
pembimbing}\label{arahan-resmi-tga-dan-gerbang-dosen-pembimbing}

\subsection{3.1 Kesiapan administratif}\label{kesiapan-administratif}

\begin{itemize}
\tightlist
\item[$\square$]
  Memastikan status mahasiswa aktif.
\item[$\square$]
  Memastikan telah lulus ujian komprehensif sesuai persyaratan pengajuan
  TGA.
\item[$\square$]
  Memastikan seluruh mata kuliah teori, studio, dan kerja praktik
  semester 1-7 sudah ditempuh, nilainya sudah keluar, dan dinyatakan
  lulus minimal D.
\item[$\square$]
  Memastikan jumlah nilai D tidak melebihi 25\% atau 36 SKS dari total
  144 SKS.
\item[$\square$]
  Memastikan bobot MK Tugas Akhir adalah 6 SKS dan tercantum dalam KRS
  sesuai ketentuan program studi.
\item[$\square$]
  Memastikan MK Tugas Akhir tercantum dalam KRS dan telah memperoleh
  persetujuan dosen pembimbing akademik.
\item[$\square$]
  Menyelesaikan pendaftaran peserta TGA melalui kanal resmi program
  studi.
\item[$\square$]
  Memastikan tidak mengambil mata kuliah teori, studio, atau kerja
  praktik bersamaan dengan MK TGA dan/atau KKN sesuai aturan yang
  berlaku.
\item[$\square$]
  Mengikuti pembekalan TGA.
\item[$\square$]
  Memeriksa laman resmi TGA untuk panduan, jadwal, formulir, rubrik,
  layanan Turnitin, dan pengumuman terbaru.
\end{itemize}

Kanal rujukan yang tercantum dalam panduan:
\texttt{https://sites.google.com/ugm.ac.id/tgapwkugm}.

\subsection{3.2 Pembimbingan dan penajaman
proposal}\label{pembimbingan-dan-penajaman-proposal}

\begin{itemize}
\tightlist
\item[$\square$]
  Menyepakati jadwal kerja TGA bersama dosen pembimbing, termasuk target
  Progres 1, Progres 2, pra-sidang, sidang, dan yudisium.
\item[$\square$]
  Mencatat tanggal Progres dari pembekalan sebagai target kerja;
  memperlakukan timeline tersebut sebagai informasi yang dapat
  disesuaikan bersama dosen pembimbing, bukan batas yang menghalangi
  penyelesaian lebih cepat.
\item[$\square$]
  Mencatat setiap pembimbingan dalam logbook TGA di SIMASTER.
\item[$\square$]
  Membawa judul kerja, rumusan masalah, tujuan, ruang lingkup, desain,
  kebutuhan data, dan risiko data ke pembimbingan awal.
\item[$\square$]
  Meminta keputusan dosen pembimbing atas struktur bab, termasuk apakah
  subbab 1.4 memang ditiadakan atau perlu ditambahkan.
\item[$\square$]
  Mengunci versi proposal yang akan digunakan sebagai dasar pengumpulan
  data.
\item[$\square$]
  Mendapatkan persetujuan proposal sebelum melakukan survei lapangan.
\item[$\square$]
  Jika diperlukan, mengajukan Surat Pengantar Survei melalui aplikasi
  persuratan setelah proposal disetujui.
\item[$\square$]
  Mencatat arahan dosen yang mengubah judul, unit, indikator, metode,
  jadwal, atau batas klaim.
\item[$\square$]
  Menghitung jumlah bimbingan dan memastikan minimal 10 kali bimbingan
  tercatat sebelum pendaftaran pra-sidang.
\end{itemize}

\subsection{3.3 Pengumpulan data, validitas data, dan Progres
1}\label{pengumpulan-data-validitas-data-dan-progres-1}

\begin{itemize}
\tightlist
\item[$\square$]
  Mengumpulkan sumber terbuka lintas Jabodetabek sesuai konstruk dan
  kemampuan akuisisi sebenarnya.
\item[$\square$]
  Mencatat versi, periode observasi, tanggal akses, unit, cakupan,
  definisi, metode pembentukan, nilai hilang, serta syarat penggunaan.
\item[$\square$]
  Membedakan observasi, agregat, estimasi, dan proksi; label D/A/S/P
  menjelaskan bentuk bukti, bukan ranking mutu.
\item[$\square$]
  Menyusun daftar cakupan pusat--indikator--tahun dan memilih seri yang
  sebanding.
\item[$\square$]
  Memeriksa duplikasi usaha, pencocokan fasilitas, atribut kapasitas
  atau akreditasi, dan dasar koefisien pekerjaan.
\item[$\square$]
  Memisahkan data kalibrasi dari pemeriksaan independen; periksa apakah
  sumber pembanding menyalin data yang sama.
\item[$\square$]
  Menyiapkan bahan Progres 1 berupa data dan audit yang benar-benar
  tersedia, beserta batasnya.
\item[$\square$]
  Mengikuti mekanisme validasi dan Progres 1 dari pembimbing; perbaiki
  modul yang belum memenuhi syarat.
\end{itemize}

\subsection{3.4 Analisis, validitas hasil, dan Progres
2}\label{analisis-validitas-hasil-dan-progres-2}

\begin{itemize}
\tightlist
\item[$\square$]
  Membekukan nilai baseline dan manifest konfigurasi bukti sebelum
  menjalankan sensitivitas.
\item[$\square$]
  Menjalankan analisis sesuai konstruk, bentuk pengukuran, hasil
  validasi, dan ketidakpastian tiap modul.
\item[$\square$]
  Menyimpan skrip, langkah transformasi, parameter, tabel antara, dan
  keluaran setiap konfigurasi.
\item[$\square$]
  Menyiapkan bahan Progres 2 yang memperlihatkan hasil analisis,
  keterbatasan, dan keputusan interpretasi.
\item[$\square$]
  Mengikuti mekanisme Progres 2 yang ditetapkan dosen pembimbing.
\item[$\square$]
  Jika hasil analisis dinyatakan tidak valid, memperbaiki analisis dan
  mengulang validasi hasil.
\item[$\square$]
  Mendapatkan persetujuan dosen pembimbing sebelum naskah diajukan ke
  pra-sidang.
\end{itemize}

\subsection{3.5 Pra-sidang}\label{pra-sidang}

\begin{itemize}
\tightlist
\item[$\square$]
  Memenuhi minimal 10 kali bimbingan yang telah diverifikasi dosen
  pembimbing.
\item[$\square$]
  Memiliki nilai Progres 1 dan Progres 2.
\item[$\square$]
  Menyelesaikan pemeriksaan isi, sitasi, tabel, gambar, bahasa, typo,
  dan konsistensi format.
\item[$\square$]
  Menjalankan pemeriksaan Turnitin dan memastikan similarity maksimal
  15\% sesuai panduan yang tersedia.
\item[$\square$]
  Mengajukan persetujuan maju pra-sidang kepada dosen pembimbing melalui
  formulir yang berlaku.
\item[$\square$]
  Mengisi pendaftaran pra-sidang melalui kesekretariatan TGA.
\item[$\square$]
  Menyiapkan PDF TGA ber-\emph{bookmark} dan PDF bahan presentasi sesuai
  instruksi pendaftaran.
\item[$\square$]
  Mengirim berkas ke alamat dan kanal yang diminta oleh kesekretariatan
  TGA setelah memeriksa kembali informasi terbaru.
\item[$\square$]
  Menyediakan hardcopy hanya jika diminta pembimbing atau penguji.
\item[$\square$]
  Menyiapkan presentasi 20 menit dan latihan jawaban untuk sesi tanya
  jawab hingga 90 menit.
\item[$\square$]
  Mengetahui tiga kemungkinan hasil: tidak lulus/mengulang, lulus dengan
  revisi, atau lulus tanpa revisi dan melanjutkan ke sidang.
\item[$\square$]
  Mencatat semua masukan pra-sidang sebagai daftar revisi yang dapat
  dilacak.
\item[$\square$]
  Jika sidang tidak dilakukan dalam satu semester setelah pra-sidang,
  memeriksa kewajiban mengulang pra-sidang dan risiko penggantian topik.
\end{itemize}

\subsection{3.6 Sidang}\label{sidang}

\begin{itemize}
\tightlist
\item[$\square$]
  Memastikan pra-sidang telah lulus dan persetujuan dosen pembimbing
  untuk mendaftar sidang telah diperoleh.
\item[$\square$]
  Memastikan minimal dua dosen penguji ditetapkan dan salah satunya
  merupakan penguji pra-sidang.
\item[$\square$]
  Mengulangi pemeriksaan akhir dokumen, typo, sitasi, Turnitin,
  \emph{bookmark}, dan PDF presentasi.
\item[$\square$]
  Mengikuti prosedur pendaftaran sidang yang sama dengan pra-sidang,
  sesuai informasi terbaru kesekretariatan.
\item[$\square$]
  Mencetak dua rangkap dokumen TGA dan menjilid \emph{hardcover} lengkap
  dengan lembar pengesahan jika diwajibkan untuk sidang.
\item[$\square$]
  Datang 30 menit sebelum sidang untuk menyiapkan file presentasi.
\item[$\square$]
  Menyiapkan presentasi 20 menit dan latihan sesi tanya jawab 90 menit.
\item[$\square$]
  Mencatat hasil sidang: lulus dan lanjut yudisium atau tidak lulus dan
  mengulang sidang.
\item[$\square$]
  Memastikan catatan dan masukan penguji dicatat oleh pembimbing,
  dipindai, dan dilampirkan di bagian belakang dokumen.
\item[$\square$]
  Memperhatikan aturan panduan yang tersedia bahwa tidak ada revisi
  pasca-sidang; seluruh revisi wajib diselesaikan sebelum dokumen final
  diajukan.
\end{itemize}

\subsection{3.7 Yudisium dan wisuda}\label{yudisium-dan-wisuda}

\begin{itemize}
\tightlist
\item[$\square$]
  Mengajukan surat bebas perpustakaan DTAP dengan PDF skripsi
  ber-\emph{bookmark} dan PDF presentasi sesuai instruksi perpustakaan.
\item[$\square$]
  Mengonfirmasi prosedur penyerahan hardfile kepada perpustakaan ketika
  diminta atau ketika kondisi memungkinkan.
\item[$\square$]
  Mengisi formulir yudisium dan \emph{exit survey}.
\item[$\square$]
  Menyelesaikan pendaftaran yudisium paling lambat 7 hari sebelum jadwal
  yudisium, kecuali aturan terbaru menetapkan lain.
\item[$\square$]
  Membuat folder Google Drive berisi foto, ijazah SMA, KTM, transkrip,
  formulir yudisium, \emph{exit survey}, skripsi final, bukti pembayaran
  UKT, KRS terakhir, dan surat pengganti lembar pengesahan bila
  diperlukan.
\item[$\square$]
  Mengirim folder kepada bagian akademik melalui prosedur yang berlaku
  dan melakukan konfirmasi.
\item[$\square$]
  Mengunggah dokumen ke ETD UGM dan melakukan konfirmasi.
\item[$\square$]
  Memastikan berita acara, SKL, transkrip, dan dokumen yudisium telah
  diproses.
\item[$\square$]
  Memeriksa jadwal wisuda setelah status yudisium selesai.
\end{itemize}

\subsection{3.8 Bobot penilaian dan risiko
keterlambatan}\label{bobot-penilaian-dan-risiko-keterlambatan}

Catatan dari panduan yang tersedia:

{\def\LTcaptype{none} % do not increment counter
\begin{longtable}[]{@{}lr@{}}
\toprule\noalign{}
Komponen & Bobot \\
\midrule\noalign{}
\endhead
\bottomrule\noalign{}
\endlastfoot
Progres 1 & 5\% \\
Progres 2 & 5\% \\
Pra-sidang & 40\%: pembimbing 50\%, penguji 50\% \\
Sidang & 50\%: pembimbing 40\%, penguji 1 30\%, penguji 2 30\% \\
\end{longtable}
}

{\def\LTcaptype{none} % do not increment counter
\begin{longtable}[]{@{}lr@{}}
\toprule\noalign{}
Keterlambatan penyelesaian TGA & Pengurangan nilai pada panduan \\
\midrule\noalign{}
\endhead
\bottomrule\noalign{}
\endlastfoot
1-2 semester & 5\% \\
3-4 semester & 10\% \\
5-6 semester & 15\% \\
7-8 semester & 20\% \\
\end{longtable}
}

\begin{itemize}
\tightlist
\item[$\square$]
  Menetapkan jadwal kerja realistis bersama dosen pembimbing dengan
  mempertimbangkan jadwal yudisium bulanan dan wisuda tiga bulanan.
\item[$\square$]
  Menyimpan tanggal aktual setiap gerbang agar risiko keterlambatan
  terlihat lebih awal.
\end{itemize}

\subsection{3.9 Rencana penggunaan 10
bimbingan}\label{rencana-penggunaan-10-bimbingan}

Jumlah minimal dan pencatatan bimbingan mengikuti panduan TGA yang
tersedia serta verifikasi kanal resmi. Rencana berikut menyediakan bahan
substantif untuk sepuluh pertemuan; Progres 1/2 tidak otomatis
menggantikan bimbingan.

{\def\LTcaptype{none} % do not increment counter
\begin{longtable}[]{@{}
  >{\raggedright\arraybackslash}p{(\linewidth - 4\tabcolsep) * \real{0.3333}}
  >{\raggedright\arraybackslash}p{(\linewidth - 4\tabcolsep) * \real{0.3333}}
  >{\raggedright\arraybackslash}p{(\linewidth - 4\tabcolsep) * \real{0.3333}}@{}}
\toprule\noalign{}
\begin{minipage}[b]{\linewidth}\raggedright
Sesi
\end{minipage} & \begin{minipage}[b]{\linewidth}\raggedright
Fokus
\end{minipage} & \begin{minipage}[b]{\linewidth}\raggedright
Bahan yang dibawa
\end{minipage} \\
\midrule\noalign{}
\endhead
\bottomrule\noalign{}
\endlastfoot
1 & Pertanyaan dan batas klaim & Bab 1, arah proksi dan sumber terbuka,
desain multitemporal sesuai bukti. \\
2 & Teori dan bukti tandingan & Kerangka Bab 2, gap terbatas korpus,
perbedaan relasi dan fungsi. \\
3 & Sumber dan unit & Inventaris sumber terbuka, cakupan
pusat--indikator--tahun, calon unit. \\
4 & Konstruksi pengukuran & Pencocokan usaha/bangunan, atribut layanan,
dasar koefisien pekerjaan. \\
5 & Relasi metropolitan & Jaringan, informasi operator, kendala estimasi
OD, rencana validasi. \\
6 & Kesebandingan waktu & Seri yang layak, perubahan pencatatan,
geometri dan kategori. \\
7 & Validitas awal & Hasil pemeriksaan independen, ketidakpastian, bahan
Progres 1. \\
8 & Fungsi relatif dan asosiasi & Tolok ukur, residual, pemeriksaan
sirkularitas, hasil Q1--Q3. \\
9 & Heterogenitas dan ketahanan & Hasil Q3, perbandingan antarpusat,
sensitivitas, konfigurasi gagal, bahan Progres 2. \\
10 & Kelayakan naskah & Bab yang selaras, batas inferensi, kesiapan
pra-sidang. \\
\end{longtable}
}

Pertemuan tambahan digunakan sesuai kebutuhan revisi dan verifikasi
administrasi, tanpa bergantung pada respons permohonan data.

\subsection{3.10 Urutan akuisisi dan keputusan
data}\label{urutan-akuisisi-dan-keputusan-data}

\begin{itemize}
\tightlist
\item[$\square$]
  Dahulukan sumber yang mencakup seluruh wilayah dengan definisi
  sebanding; periksa keunggulan cakupan, jangan mengasumsikannya.
\item[$\square$]
  Ambil sampel ekstraksi untuk memastikan lokasi, kategori, tahun, dan
  cara reproduksi sebelum pengumpulan skala kawasan.
\item[$\square$]
  Bangun modul fungsi, relasi, dan temporal secara bertahap; kegagalan
  satu modul tidak menghentikan modul lain.
\item[$\square$]
  Jika sumber tidak dapat digunakan, cari sumber terbuka lain yang
  mengukur konstruk sejenis atau nyatakan batas modul.
\item[$\square$]
  Jika data hanya mutakhir, gunakan untuk analisis periode tersebut;
  jangan menganggapnya inventaris historis.
\item[$\square$]
  Bekukan konfigurasi setelah audit dan simpan perubahan versi tanpa
  menimpa hasil sebelumnya.
\item[$\square$]
  Uji penghitungan objek dari rekaman CCTV dengan deteksi YOLO,
  pelacakan objek, garis hitung, dan sampel validasi manual; gunakan
  SUMO hanya jika data kalibrasinya memadai.
\end{itemize}

\section{4. Gerbang substantif
penelitian}\label{gerbang-substantif-penelitian}

\subsection{4.1 Penguncian pertanyaan dan
klaim}\label{penguncian-pertanyaan-dan-klaim}

\begin{itemize}
\tightlist
\item[$\square$]
  Gunakan tiga pertanyaan aktif pada Bab 1: Q1 relasi terarah; Q2 fungsi
  relatif, manfaat, beban, dan perubahan; Q3 hubungan integrasi dengan
  fungsi relatif serta konfigurasi hasil antarpusat.
\item[$\square$]
  Tempatkan sensitivitas sebagai pemeriksaan lintas pertanyaan, bukan
  pertanyaan substantif tersendiri.
\item[$\square$]
  Pertahankan perbedaan pengamatan relasi, estimasi relasi,
  aksesibilitas potensial, fungsi, kinerja, dan kesejahteraan.
\item[$\square$]
  Tetapkan konstruk, unit, dan batas interpretasi sebelum melihat hasil
  klasifikasi.
\item[$\square$]
  Jangan menyebut asosiasi atau perubahan temporal sebagai sebab-akibat
  tanpa strategi identifikasi tambahan.
\end{itemize}

\subsection{4.2 Audit literatur}\label{audit-literatur}

\begin{itemize}
\tightlist
\item[$\square$]
  Membaca teks penuh sumber prioritas sebelum menjadikan detail metode
  atau temuan sebagai dasar naskah final.
\item[$\square$]
  Mencatat definisi unit, bukti relasional, variabel fungsi/kinerja,
  ukuran lokal, metode, hasil, batasan, dan status akses setiap sumber
  kunci.
\item[$\square$]
  Menyusun Bab 2.2 sebagai perdebatan: manfaat jaringan, kemungkinan
  \emph{shadow}, masalah pengukuran, estafet Jakarta,
  akses/manfaat/beban, dan ketidakpastian metode.
\item[$\square$]
  Mengubah ringkasan artikel menjadi argumen perbandingan, bukan urutan
  73 ringkasan.
\item[$\square$]
  Menyusun tabel penelitian terdahulu dengan kolom penulis/tahun,
  fokus/lokus, metode, temuan, dan perbedaan dengan penelitian ini.
\item[$\square$]
  Menggunakan frasa ``dalam korpus yang ditelusuri'' untuk kesenjangan
  yang belum terverifikasi secara sistematis.
\item[$\square$]
  Memastikan setiap klaim rinci hanya memakai sumber yang benar-benar
  telah diperiksa dan sitasinya dapat ditelusuri.
\item[$\square$]
  Memasukkan bukti tandingan bahwa kedekatan dengan pusat besar tidak
  otomatis menghasilkan \emph{agglomeration shadow}.
\end{itemize}

\subsection{4.3 Audit sumber dan konstruksi
indikator}\label{audit-sumber-dan-konstruksi-indikator}

\begin{itemize}
\tightlist
\item[$\square$]
  Petakan sumber usaha/bangunan, fasilitas beserta atributnya,
  jaringan/operator, massa lokal, dan pengamatan multitemporal.
\item[$\square$]
  Gunakan data pemerintah terbuka apabila cocok; jangan jadikan status
  resmi sebagai standar kebenaran tunggal.
\item[$\square$]
  Periksa cakupan spasial dan historis OSM, Google Maps, direktori lain,
  serta sumber pemda pada unit yang sama.
\item[$\square$]
  Tetapkan prosedur akuisisi, deduplikasi, pencocokan lokasi,
  klasifikasi sektor, dan dokumentasi asal data.
\item[$\square$]
  Tetapkan koefisien estimasi yang memiliki dasar dan rentang pengujian;
  koefisien belum tersedia berarti estimasi belum siap, bukan izin
  mengarang angka.
\item[$\square$]
  Simpan konfigurasi bukti dengan sumber, tahun, unit, asumsi, indikator
  aktif, validasi, ketidakpastian, dan batas klaim.
\item[$\square$]
  Bandingkan sumber menurut kemampuan mengukur konstruk serta biaya dan
  keterulangan pemerolehannya.
\end{itemize}

\subsection{4.4 Isu data yang wajib
diselesaikan}\label{isu-data-yang-wajib-diselesaikan}

\begin{itemize}
\tightlist
\item[$\square$]
  Memeriksa konflik ambang WSM 2024 dalam inventaris terdahulu melalui
  dokumen publik; uji alternatif yang beralasan tanpa menunggu
  klarifikasi instansi.
\item[$\square$]
  Memastikan tahun rujukan penduduk dan versi dokumen WSM yang dipakai.
\item[$\square$]
  Menyatakan bahwa persentase orientasi WSM bukan matriks OD antarpusat.
\item[$\square$]
  Menyatakan bahwa tabel bangkitan-tarikan BPTJ bukan pasangan arus OD
  yang teramati.
\item[$\square$]
  Menyatakan bahwa jaringan dan waktu tempuh hanya menghasilkan
  aksesibilitas potensial.
\item[$\square$]
  Menyatakan bahwa POI, fasilitas, kawasan industri, morfologi
  terbangun, cahaya malam, dan direktori adalah proksi sesuai
  konstruknya masing-masing.
\item[$\square$]
  Menolak pengisian nilai hilang sebagai nol tanpa alasan substantif dan
  aturan yang terdokumentasi.
\item[$\square$]
  Menolak disagregasi nilai kabupaten/kota atau zona kasar menjadi
  variasi kecamatan/grid tanpa model eksplisit, asumsi, validasi, dan
  label sintetis.
\item[$\square$]
  Memastikan keluaran Web GIS interaktif mengikuti syarat publikasi
  sumber dan tidak memuat identitas sensitif.
\end{itemize}

\section{5. Operasionalisasi metode}\label{operasionalisasi-metode}

\subsection{5.1 Unit, pusat, dan wilayah
tangkapan}\label{unit-pusat-dan-wilayah-tangkapan}

\begin{itemize}
\tightlist
\item[$\square$]
  Mendefinisikan unit amatan sebagai satuan tempat observasi dicatat
  pada sumber.
\item[$\square$]
  Mendefinisikan unit perhitungan sebagai satuan bersama yang sah dalam
  konfigurasi bukti tertentu.
\item[$\square$]
  Mendefinisikan unit analisis sebagai pusat sekunder beserta wilayah
  tangkapan fungsionalnya.
\item[$\square$]
  Menentukan kandidat pusat dari konsentrasi aktivitas dan fungsi sumber
  terbuka, dengan dokumen perencanaan sebagai pembanding.
\item[$\square$]
  Menetapkan aturan kandidat dan \emph{catchment} sebelum menghitung
  fungsi relatif dan tipologi.
\item[$\square$]
  Memisahkan variabel delineasi dari variabel hasil sejauh data
  memungkinkan.
\item[$\square$]
  Jika indikator yang sama dipakai untuk delineasi dan hasil,
  menjalankan pemeriksaan dengan indikator tersebut dikeluarkan.
\item[$\square$]
  Menetapkan aturan agregasi desa ke kecamatan hanya jika kode,
  definisi, geometri, dan cakupan mendukung.
\item[$\square$]
  Menjalankan analisis paralel apabila unit relasional dan fungsi tidak
  dapat disetarakan.
\item[$\square$]
  Mengeluarkan modul apabila keterbandingan tidak dapat
  dipertanggungjawabkan.
\item[$\square$]
  Mencatat validasi manual tiga sampai lima pusat prioritas sebagai
  pemeriksaan purposif, bukan sampel inferensial.
\end{itemize}

\subsection{5.2 Pengukuran relasi
metropolitan}\label{pengukuran-relasi-metropolitan}

\begin{itemize}
\tightlist
\item[$\square$]
  Pisahkan pengamatan relasi, estimasi OD, dan aksesibilitas potensial.
\item[$\square$]
  Gunakan pasangan arus komuter teramati sebagai bukti relasi utama
  jika unit dan periodenya kompatibel.
\item[$\square$]
  Bangun jaringan dan biaya perjalanan dari sumber terbuka; catat
  perbedaan moda dan periode.
\item[$\square$]
  Bedakan rute, frekuensi, armada, kapasitas, jumlah penumpang, volume
  kendaraan, dan pasangan asal--tujuan.
\item[$\square$]
  Hitung intensitas arus, orientasi ke DKI, pangsa arus ke pusat
  sekunder lain, arus masuk--keluar, asimetri, dan
  \emph{self-containment} dari satu matriks yang sama apabila
  definisinya mendukung.
\item[$\square$]
  Jangan memakai proporsi menuju DKI sendirian sebagai ukuran
  integrasi.
\item[$\square$]
  Estimasikan OD hanya sebagai jalur eksploratif jika pasangan arus
  teramati tidak memadai; gunakan IPF/Furness hanya jika marginal
  kompatibel.
\item[$\square$]
  Periksa estimasi terhadap data yang tidak digunakan untuk kalibrasi.
  Kecocokan marginal tidak menjamin ketepatan pasangan OD.
\item[$\square$]
  Jangan menguji OD sintetis terhadap fungsi pekerjaan jika pekerjaan
  dipakai sebagai massa tujuan.
\item[$\square$]
  Pasangkan relasi komuter dengan pekerjaan hanya jika keduanya dibentuk
  dari sumber yang independen; domain lain tidak diberi klaim mekanisme
  khusus tanpa relasi yang sesuai.
\item[$\square$]
  Sajikan aksesibilitas potensial sebagai konstruk tersendiri; evaluasi
  CCTV/SUMO sebagai opsi setelah audit kelayakan.
\end{itemize}

\subsection{5.3 Pengukuran fungsi
relatif}\label{pengukuran-fungsi-relatif}

\begin{itemize}
\item[$\square$]
  Batasi analisis inti pada satu atau paling banyak dua domain:
  pekerjaan sebagai pilihan utama dan layanan berorde tinggi sebagai
  pilihan kedua jika lolos audit.
\item[$\square$]
  Menentukan fungsi aktual \(Y_{ik}\), ukuran lokal \(S_i\), dan
  karakteristik lokal \(X_i\) sebelum klasifikasi.
\item[$\square$]
  Gunakan seluruh pusat sekunder Jabodetabek yang kompatibel sebagai
  populasi pembanding; keluarkan DKI dari pembentukan ekspektasi.
\item[$\square$]
  Memilih model fungsi-ukuran hanya jika jumlah unit efektif, bentuk
  distribusi, pencilan, residual, dan diagnostik memadai.
\item[$\square$]
  Gunakan prediksi \emph{leave-one-out}; jika model tidak layak, gunakan
  kelompok ukuran atau tetangga ukuran terdekat yang transparan.
\item[$\square$]
  Menyatakan jika benchmark hanya relatif terhadap kawasan studi dan
  bukan ukuran universal.
\item[$\square$]
  Gunakan pembanding eksternal hanya sebagai sensitivitas tambahan dan
  jangan menjadikannya syarat analisis.
\item[$\square$]
  Menetapkan bentuk umum fungsi yang diharapkan:

  \[
  g_k(Y_{ik}) = \alpha_k + \beta_k h(S_i) + \boldsymbol{\gamma}_k^{\mathsf T}\mathbf{X}_i + \varepsilon_{ik}
  \]
\item[$\square$]
  Menghitung residual fungsi relatif:

  \[
  R_{ik}=g_k(Y_{ik})-\widehat{g_k(Y_{ik})}
  \]
\item[$\square$]
  Menafsirkan residual positif sebagai surplus relatif pada konstruk
  \(k\) dan residual negatif sebagai defisit relatif pada konstruk
  \(k\).
\item[$\square$]
  Tidak menyebut residual layanan sebagai surplus produktivitas.
\item[$\square$]
  Estimasikan pekerjaan dari usaha/bangunan dan koefisien beralasan;
  laporkan asumsi, rentang, serta validasi, bukan menyamakan jumlah
  titik dengan pekerja.
\item[$\square$]
  Tidak menyebut potret satu tahun sebagai \emph{functional upgrading}
  temporal.
\item[$\square$]
  Menetapkan aturan prediksi, ketidakpastian, transformasi, dan populasi
  pembanding sebelum klasifikasi.
\item[$\square$]
  Analisis perubahan hanya pada seri yang definisi, batas, kategori, dan
  cakupannya sebanding.
\item[$\square$]
  Bedakan perubahan fenomena dari perubahan pemetaan; jangan mengisi
  masa lalu dengan inventaris masa kini.
\item[$\square$]
  Teruskan ketidakpastian pengukuran ke residual dan asosiasi; bedakan
  acuan tetap dan pembanding per tahun.
\end{itemize}

\subsection{5.4 Manfaat, ketergantungan, dan
beban}\label{manfaat-ketergantungan-dan-beban}

\begin{itemize}
\tightlist
\item[$\square$]
  Pertahankan residu fungsi sebagai dimensi tersendiri; jangan hitung
  surplus sebagai manfaat atau defisit sebagai beban untuk kedua
  kalinya.
\item[$\square$]
  Bentuk profil manfaat dari indikator terpisah yang benar-benar
  tersedia, misalnya akses pekerjaan, penghematan waktu, atau
  upah/produktivitas jika kompatibel.
\item[$\square$]
  Bentuk profil beban dari indikator terpisah, misalnya orientasi sangat
  bergantung pada inti, durasi/biaya perjalanan, ketidakseimbangan arus,
  atau tekanan infrastruktur terukur.
\item[$\square$]
  Memeriksa apakah komuter menunjukkan akses, beban, atau keduanya.
\item[$\square$]
  Memeriksa apakah industrialisasi menunjukkan \emph{upgrading},
  produksi subordinat, atau hanya keberadaan kawasan industri.
\item[$\square$]
  Tidak menjadikan kebocoran nilai, retensi manfaat, atau
  \emph{upgrading} sebagai indikator operasional tanpa bukti yang
  sesuai.
\item[$\square$]
  Menentukan polaritas per konstruk, bukan dari tanda positif atau
  negatif mentah.
\item[$\square$]
  Mempertahankan manfaat dan beban sebagai dimensi terpisah sebelum
  membuat tipologi.
\end{itemize}

\subsection{5.5 Tipologi dan empat gerbang
label}\label{tipologi-dan-empat-gerbang-label}

\begin{itemize}
\item[$\square$]
  Menyajikan vektor profil multidimensi sebelum membuat kelas ringkas.
\item[$\square$]
  Jika skor ringkas digunakan, menetapkan indikator dan bobot dalam satu
  konfigurasi bukti yang tetap:

  \[
  Z_{id}=\frac{\sum_{k \in K_d} w_k z_{ik}}{\sum_{k \in K_d} w_k}, \qquad w_k \geq 0
  \]
\item[$\square$]
  Tidak mengubah nilai hilang menjadi nol.
\item[$\square$]
  Membekukan populasi pembanding dan normalisasi baseline sebelum
  sensitivitas.
\item[$\square$]
  Menentukan ambang baseline dari dasar teoretis, empiris, atau
  kebijakan setelah distribusi data diperiksa.
\item[$\square$]
  Menguji ambang tersebut sebagai sensitivitas, bukan mengubahnya agar
  label sesuai harapan.
\item[$\square$]
  Mempertahankan kelas relatif independen, pola konsisten dengan
  \emph{borrowed size}, campuran, pola konsisten dengan
  \emph{agglomeration shadow}, terintegrasi tetapi tertinggal, periferi
  lemah, dan tidak terselesaikan sesuai bukti.
\item[$\square$]
  Menetapkan bahwa ``relatif independen'' bukan bukti autarki.
\item[$\square$]
  Gunakan label \emph{borrowed size} hanya jika relasi tinggi, residu
  positif, dan beban independen tidak tinggi.
\item[$\square$]
  Menetapkan bahwa kelas campuran dapat memuat manfaat dan beban yang
  sama-sama tinggi atau residual yang berbeda antarkonstruk.
\item[$\square$]
  Gunakan istilah \emph{transisional} hanya untuk perubahan antarperiode
  yang sebanding; ketidakstabilan spesifikasi masuk kelas tidak
  terselesaikan.
\item[$\square$]
  Gunakan label \emph{agglomeration shadow} hanya jika integrasi tinggi,
  residu negatif, dan beban atau ketergantungan independen terbukti;
  jika bukti terakhir tidak ada, gunakan kelas terintegrasi tetapi
  tertinggal.
\item[$\square$]
  Menetapkan bahwa kelas tidak terselesaikan adalah keluaran sah ketika
  komponen wajib hilang, unit tidak kompatibel, atau hasil tidak stabil.
\end{itemize}

Label ketat hanya boleh digunakan jika seluruh gerbang berikut
terlewati:

\begin{itemize}
\tightlist
\item[$\square$]
  \textbf{Gerbang 1 - Bukti relasional:} hubungan dengan Jakarta atau
  jaringan pusat didukung pengamatan atau estimasi yang diperiksa
  independen; kedekatan saja belum cukup.
\item[$\square$]
  \textbf{Gerbang 2 - Pembanding substantif:} fungsi aktual dibandingkan
  dengan ekspektasi berdasarkan ukuran dan karakteristik lokal pada unit
  yang memadai.
\item[$\square$]
  \textbf{Gerbang 3 - Kompatibilitas:} unit, wilayah tangkapan, periode,
  populasi sasaran, dan konstruk integrasi-fungsi dapat dihubungkan
  secara sah.
\item[$\square$]
  \textbf{Gerbang 4 - Ketahanan:} label tidak runtuh pada sensitivitas
  utama dan tidak bergantung pada satu proksi bermasalah.
\end{itemize}

Jika satu gerbang gagal, gunakan ``profil'', ``surplus/defisit fungsi
relatif'', ``orientasi'', ``aksesibilitas potensial'', ``ketergantungan
komuter'', atau ``pola yang konsisten dengan'', sesuai bukti yang
tersedia.

\subsection{5.6 Sensitivitas dan
skenario}\label{sensitivitas-dan-skenario}

\begin{itemize}
\tightlist
\item[$\square$]
  Menetapkan sensitivitas indikator dan pemeriksaan korelasi/duplikasi.
\item[$\square$]
  Jalankan \emph{leave-one-indicator-out} hanya jika satu dimensi
  benar-benar memuat lebih dari satu indikator.
\item[$\square$]
  Uji satu spesifikasi alternatif tolok ukur internal sebagai
  pemeriksaan minimum.
\item[$\square$]
  Menguji bobot alternatif hanya jika skor komposit digunakan.
\item[$\square$]
  Menguji ambang alternatif.
\item[$\square$]
  Menguji unit, agregasi, atau delineasi alternatif hanya jika alternatif
  tersebut sudah tersedia dari pengolahan utama dan operasinya sah.
\item[$\square$]
  Menguji definisi pusat dan \emph{catchment} alternatif jika lebih dari
  satu definisi dapat dipertanggungjawabkan.
\item[$\square$]
  Menguji perubahan status bukti sebagai konfigurasi bukti terpisah.
\item[$\square$]
  Menggunakan autokorelasi spasial atau regresi spasial hanya jika
  jumlah unit, geometri, matriks bobot, dan asumsi model memadai.
\item[$\square$]
  Menampilkan peta perpindahan kelas, rentang skor, unit tidak
  terselesaikan, dan alasan perubahan.
\item[$\square$]
  Tidak meringkas ketahanan hanya sebagai persentase kelas yang sama.
\item[$\square$]
  Jangan membuat dataset atau model baru semata-mata untuk memperbanyak
  analisis sensitivitas.
\end{itemize}

Pisahkan tiga jenis perubahan:

\begin{itemize}
\tightlist
\item[$\square$]
  \textbf{Skenario bukti (\texttt{EV-*}):} sumber, tahun, unit, status
  mutu, dan indikator aktif berubah; perbedaan hasil berarti kesimpulan
  bergantung pada konfigurasi bukti.
\item[$\square$]
  \textbf{Skenario parameter (\texttt{P-*}):} bobot, ambang,
  normalisasi, indikator, atau spesifikasi berubah dalam konfigurasi
  bukti yang sama; ini adalah sensitivitas model.
\item[$\square$]
  \textbf{Skenario nilai indikator:} nilai indikator diubah dalam
  rentang terkalibrasi untuk pertanyaan \emph{what-if}; ini bukan
  observasi, bukan diagnosis \emph{status quo}, dan bukan prediksi
  kausal.
\end{itemize}

\section{6. Checklist penulisan skripsi sesuai template
TGA}\label{checklist-penulisan-skripsi-sesuai-template-tga}

Template menyatakan bahwa judul bab, subbab, dan konten dapat
disesuaikan dengan kebutuhan serta arahan dosen pembimbing. Struktur di
bawah ini adalah struktur kerja, bukan pengganti keputusan pembimbing.

\subsection{6.1 Bagian awal}\label{bagian-awal}

\begin{itemize}
\tightlist
\item[$\square$]
  Menetapkan judul Indonesia dan judul Inggris setelah judul kerja
  disetujui.
\item[$\square$]
  Menyusun halaman judul dengan nama, NIM, program studi, departemen,
  fakultas, universitas, dan tahun.
\item[$\square$]
  Menyusun halaman pernyataan keaslian dan kepatuhan etika.
\item[$\square$]
  Menyusun kata pengantar maksimal dua halaman.
\item[$\square$]
  Menulis intisari maksimal 250 kata yang memuat urgensi, tujuan,
  metode, dan temuan; tulis setelah temuan final tersedia.
\item[$\square$]
  Menulis abstrak bahasa Inggris maksimal 250 kata dengan isi yang
  setara.
\item[$\square$]
  Menjelaskan istilah teknis dalam bahasa yang dapat dipahami pembaca
  lintas disiplin.
\item[$\square$]
  Menetapkan maksimal lima kata atau frasa kunci untuk setiap bahasa.
\item[$\square$]
  Membuat daftar isi, daftar tabel, dan daftar gambar dengan penomoran
  otomatis.
\end{itemize}

\subsection{6.2 Bab 1 - Pendahuluan}\label{bab-1---pendahuluan}

\begin{itemize}
\tightlist
\item[$\square$]
  Menulis \texttt{1.1\ Latar\ Belakang} maksimal dua halaman.
\item[$\square$]
  Menjelaskan fenomena metropolitan Jakarta dan masalah penguatan fungsi
  lokal versus ketergantungan.
\item[$\square$]
  Menyajikan temuan penelitian terdahulu secara singkat dalam dua atau
  tiga paragraf.
\item[$\square$]
  Menjelaskan gap dan perbedaan penelitian tanpa klaim ``pertama'' yang
  belum terverifikasi.
\item[$\square$]
  Menulis \texttt{1.2\ Rumusan\ Masalah} sebagai pertanyaan yang jelas
  dan diberi pengantar singkat.
\item[$\square$]
  Menulis \texttt{1.3\ Tujuan\ Penelitian} yang langsung menjawab
  rumusan masalah.
\item[$\square$]
  Mengonfirmasi kepada dosen pembimbing apakah \texttt{1.4} memang tidak
  digunakan dalam template atau perlu diisi.
\item[$\square$]
  Menulis \texttt{1.5\ Metode\ Penelitian} sebagai ringkasan Bab 3
  maksimal sekitar 500 kata.
\item[$\square$]
  Menulis \texttt{1.6\ Manfaat\ Penelitian} dengan manfaat teoretis dan
  praktis yang benar-benar dapat ditopang temuan.
\item[$\square$]
  Menulis \texttt{1.7\ Struktur\ Penulisan} yang memperlihatkan benang
  merah Bab 1 sampai Bab 6.
\item[$\square$]
  Memastikan istilah ``integrasi'', ``kematangan'', ``borrowed size'',
  dan ``agglomeration shadow'' konsisten dengan Bab 2 dan Bab 3.
\end{itemize}

\subsection{6.3 Bab 2 - Kajian Pustaka}\label{bab-2---kajian-pustaka}

\begin{itemize}
\tightlist
\item[$\square$]
  Menulis \texttt{2.1\ Konsep-konsep\ Kunci} dari definisi kerja yang
  dipilih, bukan tumpukan definisi.
\item[$\square$]
  Menjelaskan pusat sekunder, integrasi fungsional, fungsi/kinerja,
  kematangan lokal, \emph{borrowed size}, \emph{agglomeration shadow},
  kondisi campuran, geografi fungsional, dan skenario sesuai kebutuhan.
\item[$\square$]
  Menulis \texttt{2.2\ Kajian\ Penelitian\ Terdahulu} sebagai sintesis
  perdebatan dan estafet penelitian.
\item[$\square$]
  Membandingkan fokus, lokus, unit, data, metode, temuan, keterbatasan,
  dan perbedaan dengan penelitian ini.
\item[$\square$]
  Menjelaskan bahwa polisentrisitas morfologis tidak otomatis
  membuktikan integrasi fungsional.
\item[$\square$]
  Menjelaskan bahwa aksesibilitas potensial tidak sama dengan arus
  aktual.
\item[$\square$]
  Menjelaskan bukti lokal Jakarta mengenai pascasuburbanisasi,
  dekonsentrasi pekerjaan, komuter, akses, ketimpangan, dan tata kelola.
\item[$\square$]
  Menyusun gap defensif berdasarkan kombinasi konstruk, unit, bukti
  relasional, manfaat/beban, dan ketahanan hasil.
\item[$\square$]
  Menulis \texttt{2.3\ Kerangka\ Teori/Kerangka\ Konseptual} dari empat
  lapisan: ekonomi aglomerasi, polarisasi/kausasi kumulatif, sistem
  perkotaan/jaringan, dan diagnosis \emph{borrowed size-shadow}.
\item[$\square$]
  Menyusun rantai mekanisme dari ukuran lokal, jaringan, fungsi yang
  diharapkan, fungsi aktual, residual, manfaat, beban, dan tipologi.
\item[$\square$]
  Menetapkan proposisi sebagai ekspektasi interpretif jika desain tidak
  mendukung hipotesis inferensial.
\item[$\square$]
  Menyajikan diagram kerangka konseptual buatan sendiri dan
  menyebutkannya dalam teks.
\item[$\square$]
  Menyertakan rumus atau model hanya setelah arti setiap konstruk dan
  batas penggunaannya jelas.
\item[$\square$]
  Memastikan seluruh sitasi dalam Bab 2 dapat ditelusuri ke bacaan yang
  benar-benar diperiksa.
\end{itemize}

\subsection{6.4 Bab 3 - Metode
Penelitian}\label{bab-3---metode-penelitian}

\begin{itemize}
\tightlist
\item[$\square$]
  Mempertahankan enam subbab utama template: \texttt{3.1\ Pendekatan},
  \texttt{3.2\ Desain}, \texttt{3.3\ Ruang\ Lingkup},
  \texttt{3.4\ Unit\ Amatan\ dan\ Unit\ Analisis},
  \texttt{3.5\ Metode\ Pengumpulan\ Data}, dan
  \texttt{3.6\ Metode\ Analisis\ Data}.
\item[$\square$]
  Menulis \texttt{3.1} sebagai pendekatan kuantitatif/deduktif dengan
  analisis spasial-relasional.
\item[$\square$]
  Menjelaskan mengapa pendekatan tersebut sesuai dengan Q1-Q3.
\item[$\square$]
  Menyatakan bahwa audit dokumen dan provenance adalah prosedur mutu
  data, bukan alasan untuk menyebut penelitian campuran.
\item[$\square$]
  Menulis \texttt{3.2} sebagai studi kasus tunggal kuantitatif dengan
  rancangan minimum satu periode, satu relasi komuter, satu domain
  fungsi independen, dan tolok ukur internal.
\item[$\square$]
  Menjelaskan kelebihan, keterbatasan, dan cara pengelolaan keterbatasan
  tanpa menghapusnya.
\item[$\square$]
  Menulis \texttt{3.3} dengan batas spasial Jabodetabek, pembanding
  internal sebagai garis dasar, WSM eksternal hanya sebagai
  sensitivitas, dan maksimal dua domain inti.
\item[$\square$]
  Menjelaskan bahwa jejak fungsional bukan perluasan yurisdiksi Jakarta.
\item[$\square$]
  Menulis \texttt{3.4} dengan pemisahan unit amatan, unit perhitungan,
  unit analisis, pusat sekunder, \emph{catchment}, dan lapisan
  administrasi.
\item[$\square$]
  Menulis \texttt{3.5} sebagai akuisisi sumber terbuka, konstruksi
  proksi, estimasi, dan pemeriksaan independen.
\item[$\square$]
  Menjelaskan instrumen sekunder: log akses, formulir audit, kamus data,
  crosswalk, manifest, lembar validasi manual, dan skrip reproduksi.
\item[$\square$]
  Menjelaskan populasi/cakupan, aturan sampel jika ada, serta status
  validasi manual sebagai purposif.
\item[$\square$]
  Menjelaskan bentuk bukti D/A/S/P/\(\varnothing\), asumsi estimasi,
  pemeriksaan independen, lisensi, serta privasi.
\item[$\square$]
  Menulis \texttt{3.6} secara berurutan dari audit dan pembekuan
  baseline sampai keluaran inti; Web GIS interaktif ditempatkan sebagai
  keluaran tambahan.
\item[$\square$]
  Menyertakan aturan hitung integrasi, fungsi relatif, residual, skor
  ringkas jika dipakai, nilai hilang, ambang, klasifikasi, dan
  sensitivitas.
\item[$\square$]
  Menyertakan matriks keterlacakan Q1-Q3: pertanyaan, konstruk, bukti
  minimum, unit, operasi, keluaran, dan batas klaim.
\item[$\square$]
  Menempatkan Web GIS interaktif sebagai implementasi tambahan yang
  dapat ditiadakan, bukan sumber klaim empiris baru.
\item[$\square$]
  Mengintegrasikan validitas, reliabilitas, reproduktibilitas, dan
  aturan keputusan ke bagian metode yang relevan.
\item[$\square$]
  Menghapus pengulangan memorandum, checklist internal, dan catatan
  untuk model dari naskah Bab 3 final.
\item[$\square$]
  Menambahkan sitasi metode hanya setelah sumber metodenya benar-benar
  diverifikasi.
\end{itemize}

\subsection{6.5 Bab 4 - Gambaran Umum Lokasi
Penelitian}\label{bab-4---gambaran-umum-lokasi-penelitian}

\begin{itemize}
\tightlist
\item[$\square$]
  Menjelaskan karakteristik Jabodetabek yang relevan dengan pertanyaan
  penelitian, bukan menyalin profil wilayah.
\item[$\square$]
  Menjelaskan alasan Jabodetabek dipilih sebagai kasus.
\item[$\square$]
  Menjelaskan relasi DKI sebagai inti, pusat sekunder, wilayah
  tangkapan, dan batas administratif.
\item[$\square$]
  Menampilkan peta batas administrasi tetap dan sumber geometri.
\item[$\square$]
  Menampilkan struktur penduduk, jaringan, pekerjaan, layanan, dan
  aktivitas hanya sejauh relevan dengan analisis.
\item[$\square$]
  Menjelaskan kandidat pusat sekunder dan dasar awal pemilihannya tanpa
  menganggap status administratif sebagai bukti tunggal.
\item[$\square$]
  Memisahkan informasi konteks Bab 4 dari temuan diagnosis yang baru
  boleh muncul di Bab 5.
\end{itemize}

\subsection{6.6 Bab 5 - Temuan dan
Pembahasan}\label{bab-5---temuan-dan-pembahasan}

\begin{itemize}
\tightlist
\item[$\square$]
  Menyusun temuan berdasarkan Q1-Q3, bukan berdasarkan urutan dataset.
\item[$\square$]
  Menampilkan tabel audit dan manifest konfigurasi bukti yang digunakan
  sebagai dasar hasil.
\item[$\square$]
  Menampilkan profil/peta integrasi atau aksesibilitas sesuai tingkat M.
\item[$\square$]
  Menampilkan fungsi aktual, fungsi yang diharapkan, dan residual per
  konstruk sesuai tingkat U.
\item[$\square$]
  Menampilkan manfaat, beban, ketergantungan, dan kelas campuran secara
  terpisah.
\item[$\square$]
  Menampilkan gradien dan tipologi \emph{status quo} dengan batas klaim
  yang terlihat.
\item[$\square$]
  Menampilkan hasil sensitivitas, perpindahan kelas, unit tidak
  terselesaikan, dan alasan perubahan.
\item[$\square$]
  Menjelaskan perbedaan skenario bukti, skenario parameter, dan skenario
  nilai indikator.
\item[$\square$]
  Menulis pembahasan dengan membandingkan temuan terhadap teori dan
  penelitian terdahulu.
\item[$\square$]
  Menjelaskan implikasi perencanaan menurut tanggung jawab administratif
  tanpa mengusulkan perubahan batas secara otomatis.
\item[$\square$]
  Tidak menyebut hubungan sebagai kausal, temporal, atau prediktif
  melampaui desain multitemporal sesuai bukti.
\item[$\square$]
  Tidak mengubah hasil Web GIS menjadi prediksi kebijakan atau ramalan
  masa depan.
\end{itemize}

\subsection{6.7 Bab 6 - Kesimpulan dan
Rekomendasi}\label{bab-6---kesimpulan-dan-rekomendasi}

\begin{itemize}
\tightlist
\item[$\square$]
  Menjawab Q1-Q3 secara ringkas dan parafrastik.
\item[$\square$]
  Membedakan temuan yang langsung didukung data dari interpretasi yang
  bersyarat.
\item[$\square$]
  Menyatakan apakah hasil mendukung pola manfaat, \emph{shadow},
  campuran, periferi lemah, atau tidak terselesaikan pada konstruk
  tertentu.
\item[$\square$]
  Menjelaskan kontribusi teoretis sesuai gap yang benar-benar terisi.
\item[$\square$]
  Menulis rekomendasi kebijakan yang sesuai dengan tingkat bukti dan
  kewenangan pemerintah yang relevan.
\item[$\square$]
  Tidak mengubah jejak fungsional menjadi rekomendasi perubahan batas
  administratif tanpa analisis kelembagaan tambahan.
\item[$\square$]
  Menulis keterbatasan desain, data, unit, proksi, waktu, dan
  generalisasi secara jujur.
\item[$\square$]
  Menyusun peluang riset lanjutan untuk perluasan seri waktu yang
  sebanding, validasi OD, kinerja ekonomi, kapasitas kelembagaan, atau
  desain kausal.
\end{itemize}

\subsection{6.8 Daftar pustaka dan
lampiran}\label{daftar-pustaka-dan-lampiran}

\begin{itemize}
\tightlist
\item[$\square$]
  Memasukkan hanya sumber yang disitasi dalam teks.
\item[$\square$]
  Menggunakan gaya sitasi penulis-tahun secara konsisten dan
  menyelaraskannya dengan daftar pustaka.
\item[$\square$]
  Menambahkan nomor halaman pada kutipan langsung dan memastikan format
  kutipan konsisten.
\item[$\square$]
  Memeriksa DOI, judul, tahun, volume, nomor, halaman, dan jenis sumber.
\item[$\square$]
  Menyimpan tabel audit data, kamus data, log akses, crosswalk,
  manifest, parameter, dan daftar data yang dikeluarkan sebagai lampiran
  jika sesuai dengan lisensi.
\item[$\square$]
  Menyimpan dokumen administrasi TGA yang relevan dan log sumber
  terbuka; paket permohonan data lama tidak dipakai.
\item[$\square$]
  Tidak melampirkan data terbatas yang melanggar perjanjian penggunaan
  atau memungkinkan rekonstruksi data sensitif.
\end{itemize}

\section{7. Pembuatan versi LaTeX
printable}\label{pembuatan-versi-latex-printable}

Markdown adalah checklist utama dan sumber isi. LaTeX/PDF hanya
representasi tata letak formal yang lebih nyaman untuk dibaca, dibuka
berulang, atau dicetak; ia bukan naskah skripsi baru dan tidak
menggantikan berkas Bab 1-3.

\subsection{7.1 Persiapan format
printable}\label{persiapan-format-printable}

\begin{itemize}
\tightlist
\item[$\square$]
  Mengonfirmasi judul roadmap dan tanggal pembaruan.
\item[$\square$]
  Menetapkan ukuran A4, margin, ukuran huruf, spasi, dan gaya heading
  yang nyaman untuk dibaca berulang.
\item[$\square$]
  Menetapkan header/footer, nomor halaman, daftar isi, dan bookmark PDF.
\item[$\square$]
  Memastikan kotak checklist terlihat jelas dan tetap dapat dicentang
  setelah dicetak.
\item[$\square$]
  Memastikan rumus, simbol, istilah miring, URL, dan nama berkas dapat
  dirender tanpa kehilangan makna.
\end{itemize}

\subsection{7.2 Sinkronisasi Markdown ke LaTeX dan
PDF}\label{sinkronisasi-markdown-ke-latex-dan-pdf}

\begin{itemize}
\tightlist
\item[$\square$]
  Mengonversi isi Markdown ke sumber LaTeX tanpa mengubah substansi
  checklist.
\item[$\square$]
  Mempertahankan seluruh heading, paragraf, daftar, tabel, rumus, dan
  log pembimbingan.
\item[$\square$]
  Tidak menjadikan sumber LaTeX sebagai tempat utama untuk mengedit isi;
  perubahan substantif dilakukan terlebih dahulu di Markdown.
\item[$\square$]
  Membuat ulang sumber LaTeX dan PDF setelah Markdown berubah secara
  material.
\item[$\square$]
  Memastikan versi LaTeX dan PDF memiliki tanggal/isi yang sama dengan
  Markdown.
\end{itemize}

\subsection{7.3 Pemeriksaan LaTeX dan
PDF}\label{pemeriksaan-latex-dan-pdf}

\begin{itemize}
\tightlist
\item[$\square$]
  Memastikan PDF berhasil dibuka, berukuran A4, dan seluruh halaman
  memiliki nomor halaman yang benar.
\item[$\square$]
  Memastikan daftar isi dan bookmark menuju bagian yang tepat.
\item[$\square$]
  Memastikan semua tabel, rumus, checkbox, URL, dan karakter khusus
  terbaca.
\item[$\square$]
  Memastikan tidak ada halaman kosong, teks keluar margin, tabel
  terpotong, atau baris yang bertumpuk.
\item[$\square$]
  Memastikan tabel log bimbingan memiliki ruang yang cukup untuk diisi
  atau diberi catatan.
\item[$\square$]
  Membandingkan jumlah dan urutan checklist PDF dengan Markdown.
\item[$\square$]
  Melakukan pemeriksaan visual terakhir pada layar dan hasil cetak jika
  diperlukan.
\end{itemize}

\section{8. Artefak wajib per gerbang}\label{artefak-wajib-per-gerbang}

{\def\LTcaptype{none} % do not increment counter
\begin{longtable}[]{@{}
  >{\raggedright\arraybackslash}p{(\linewidth - 4\tabcolsep) * \real{0.3333}}
  >{\raggedright\arraybackslash}p{(\linewidth - 4\tabcolsep) * \real{0.3333}}
  >{\raggedright\arraybackslash}p{(\linewidth - 4\tabcolsep) * \real{0.3333}}@{}}
\toprule\noalign{}
\begin{minipage}[b]{\linewidth}\raggedright
Gerbang
\end{minipage} & \begin{minipage}[b]{\linewidth}\raggedright
Artefak minimum
\end{minipage} & \begin{minipage}[b]{\linewidth}\raggedright
Persetujuan/cek
\end{minipage} \\
\midrule\noalign{}
\endhead
\bottomrule\noalign{}
\endlastfoot
Proposal & judul, Q1-Q3, tujuan, ruang lingkup, desain, kebutuhan data,
batas klaim & dosen pembimbing \\
Data & log akses, registri provenance, kamus data, geometri/crosswalk,
audit mutu & validitas data/Progres 1 \\
Baseline & nilai baseline beku, manifest, unit, tahun, indikator, status
D/A/S/P/\(\varnothing\) & keputusan metode \\
Analisis & tabel hasil, skrip/langkah transformasi, residual, profil
manfaat-beban, kelas/gradien & validitas hasil/Progres 2 \\
Sensitivitas & konfigurasi \texttt{EV-*}, \texttt{P-*}, hasil
perpindahan kelas, unit tidak terselesaikan & dosen pembimbing \\
Naskah & Bab 1-6, daftar pustaka, lampiran, tabel/gambar, PDF
ber-\emph{bookmark} & persetujuan pra-sidang \\
Pra-sidang & pendaftaran, logbook minimal 10 bimbingan, nilai Progres
1/2, Turnitin & kesekretariatan TGA \\
Sidang & naskah final, PDF presentasi, dua \emph{hardcover} jika
diwajibkan & DPS dan penguji \\
Yudisium & surat bebas perpustakaan, folder dokumen, formulir,
\emph{exit survey}, ETD & bagian akademik/ETD \\
\end{longtable}
}

\section{9. Pemeriksaan akhir sebelum mencentang
selesai}\label{pemeriksaan-akhir-sebelum-mencentang-selesai}

\begin{itemize}
\tightlist
\item[$\square$]
  Tidak ada data yang disebut diterima padahal baru direncanakan.
\item[$\square$]
  Tidak ada proksi yang disebut sebagai pengamatan langsung.
\item[$\square$]
  Tidak ada data agregat yang diturunkan secara semu ke unit lebih
  rinci.
\item[$\square$]
  Tidak ada label \emph{borrowed size} atau \emph{agglomeration shadow}
  tanpa empat gerbang bukti.
\item[$\square$]
  Tidak ada perubahan skenario yang disajikan sebagai perubahan kondisi
  metropolitan.
\item[$\square$]
  Tidak ada klaim kausal, temporal, prediktif, atau perubahan yurisdiksi
  yang melampaui desain.
\item[$\square$]
  Tahun sumber, unit, status mutu, dan batas klaim terlihat pada
  tabel/peta utama.
\item[$\square$]
  Hasil campuran dan ketidakstabilan tidak disembunyikan demi satu label
  tunggal.
\item[$\square$]
  Semua keputusan yang masih bersyarat ditandai dan dibawa ke
  pembimbingan.
\item[$\square$]
  Naskah final mengikuti template dan arahan dosen pembimbing, bukan
  hanya struktur checklist ini.
\end{itemize}

\section{10. Log pembimbingan dan
keputusan}\label{log-pembimbingan-dan-keputusan}

Gunakan log ini untuk menjaga agar perubahan dari pembimbingan tidak
hilang dari naskah, data, atau konfigurasi analisis. Logbook resmi tetap
mengikuti kanal TGA/SIMASTER.

\subsection{10.1 Log bimbingan}\label{log-bimbingan}

{\def\LTcaptype{none} % do not increment counter
\begin{longtable}[]{@{}
  >{\raggedright\arraybackslash}p{(\linewidth - 12\tabcolsep) * \real{0.1429}}
  >{\raggedright\arraybackslash}p{(\linewidth - 12\tabcolsep) * \real{0.1429}}
  >{\raggedright\arraybackslash}p{(\linewidth - 12\tabcolsep) * \real{0.1429}}
  >{\raggedright\arraybackslash}p{(\linewidth - 12\tabcolsep) * \real{0.1429}}
  >{\raggedright\arraybackslash}p{(\linewidth - 12\tabcolsep) * \real{0.1429}}
  >{\raggedright\arraybackslash}p{(\linewidth - 12\tabcolsep) * \real{0.1429}}
  >{\raggedright\arraybackslash}p{(\linewidth - 12\tabcolsep) * \real{0.1429}}@{}}
\toprule\noalign{}
\begin{minipage}[b]{\linewidth}\raggedright
Ke
\end{minipage} & \begin{minipage}[b]{\linewidth}\raggedright
Tanggal
\end{minipage} & \begin{minipage}[b]{\linewidth}\raggedright
Agenda/bahan
\end{minipage} & \begin{minipage}[b]{\linewidth}\raggedright
Arahan atau keputusan dosen
\end{minipage} & \begin{minipage}[b]{\linewidth}\raggedright
Dampak pada naskah/data
\end{minipage} & \begin{minipage}[b]{\linewidth}\raggedright
Tindak lanjut
\end{minipage} & \begin{minipage}[b]{\linewidth}\raggedright
Bukti/tautan
\end{minipage} \\
\midrule\noalign{}
\endhead
\bottomrule\noalign{}
\endlastfoot
1 & & 18 September 2026 & Proksi dan sumber terbuka sebagai fondasi;
multitemporal sesuai bukti & Keputusan pengguna & Bab 1--3, Argument,
roadmap, panduan & Ketersediaan dan mutu per modul \\
2 & & & & & & \\
3 & & & & & & \\
4 & & & & & & \\
5 & & & & & & \\
6 & & & & & & \\
7 & & & & & & \\
8 & & & & & & \\
9 & & & & & & \\
10 & & & & & & \\
11 & & & & & & \\
12 & & & & & & \\
\end{longtable}
}

\subsection{10.2 Log keputusan
substantif}\label{log-keputusan-substantif}

{\def\LTcaptype{none} % do not increment counter
\begin{longtable}[]{@{}
  >{\raggedright\arraybackslash}p{(\linewidth - 8\tabcolsep) * \real{0.2000}}
  >{\raggedright\arraybackslash}p{(\linewidth - 8\tabcolsep) * \real{0.2000}}
  >{\raggedright\arraybackslash}p{(\linewidth - 8\tabcolsep) * \real{0.2000}}
  >{\raggedright\arraybackslash}p{(\linewidth - 8\tabcolsep) * \real{0.2000}}
  >{\raggedright\arraybackslash}p{(\linewidth - 8\tabcolsep) * \real{0.2000}}@{}}
\toprule\noalign{}
\begin{minipage}[b]{\linewidth}\raggedright
Tanggal
\end{minipage} & \begin{minipage}[b]{\linewidth}\raggedright
Keputusan yang dikunci
\end{minipage} & \begin{minipage}[b]{\linewidth}\raggedright
Dasar atau alasan
\end{minipage} & \begin{minipage}[b]{\linewidth}\raggedright
Bagian yang terdampak
\end{minipage} & \begin{minipage}[b]{\linewidth}\raggedright
Perlu divalidasi ulang?
\end{minipage} \\
\midrule\noalign{}
\endhead
\bottomrule\noalign{}
\endlastfoot
& & & & \\
& & & & \\
& & & & \\
\end{longtable}
}

\end{document}
